% ---------------------------------------------------------------------------
% Experiment History -- subsection 6: Untrainability of Policy-Gradient RL
% over Raw IQ.  Sources: README A.5 (sim06 jammer, 06b, 07) + objective
% table in README A.0.  Jobs: 101622/101632/101650/101657/101817.
%
% Rewritten 2026-09-11: cut from ~700 words to ~330, to the register of the
% user-written parts 2 and 4 (short sentences, plain vocabulary, numbers
% strung into one sentence, a contrastive closer).  The tables carry the
% numbers so the prose does not have to.
%
% Two ordering decisions are deliberate and survive the cut:
%   (1) leads with the ACTION PARAMETERISATION, not the reward -- a reader
%       who takes away "they should have tuned beta" took away the wrong
%       thing;
%   (2) the four-subcarrier "cliff" is an OBSERVATION here, explained away
%       in part 7.  Do not promote it to a finding.
%
% Figure: artifacts/sim06b/jammer/run001.png, CROPPED TO THE BOTTOM TWO
% PANELS (reward falling, entropy rising).  The top panels dilute it.
% ---------------------------------------------------------------------------

\subsection{Untrainability of Policy-Gradient RL over Raw IQ}
\label{subsec:exphist:untrainable}
Direct gradients worked only because the jammer could differentiate the detector and see the symbol it
was attacking.
A frozen CNN read through its output gives neither, so the attacker has to be trained by policy
gradient, and the question was whether that works at all.

We kept the OFDM link and the frozen detector and changed only the training rule: two cooperative flows
trained with MAPPO, each injecting a complex vector per subcarrier, with $D = -\log(1-\hat{p}+\epsilon)$.
All three runs failed the same way, with $\hat{p}$ pinned at $0.999$ from start to finish.
Warming the penalty let power diverge from 4 to 33, reshaping it changed its size but not its
variation, and removing the entropy bonus froze the policy for 290 iterations.

A static probe explains the pinning: one jammed subcarrier goes unflagged at power 16
($\hat{p} \approx 0.007$), while all 64 jammed at power $0.01$ each are flagged at $1.000$.
A flow initialised near $\mathcal{N}(0,\mathbf{I})$ over 128 dimensions is broadband noise, so every
action the agent ever sampled was flagged and the reward told it nothing.
The penalty shape makes that permanent: a log barrier at $\hat{p} \approx 0.999$ is huge but nearly the
same for every element of the batch, and PPO normalises advantages within the batch, so a huge constant
penalty and no penalty are the same gradient.

Cutting the action space to two real dimensions separated the two failures.
Stealth was then solved ($\hat{p} \approx 0.003$, a genuinely varying signal) and waveform learning
still failed, with per-subcarrier BER stuck at $0.35$ against an attainable $1.0$ and no correlation
between the observed symbol and the emitted one.
A scalar frame-level reward cannot teach an input--output correlation, because the agent has to sample
the right input-dependent direction by chance before it can reinforce it.

\begin{figure}[!t]
\centering
% \includegraphics[width=\columnwidth]{fig/sim06b_run001_crop.pdf}
\caption{Two real action dimensions against the frozen detector. Mean reward falls while policy entropy
rises: the policy is diffusing, not converging.}
\label{fig:diffusion}
\end{figure}

The last variant removed the genie, giving the jammer the previous symbol instead of the current one.
For i.i.d.\ QPSK that says nothing about the symbol under attack, so cancellation is unreachable rather
than merely hard.
Five runs each repaired the previous one's mechanical failure and each failed identically
(Table~\ref{tab:blackbox}), with entropy flat at $178.4$ throughout.

\begin{table}[!t]
\renewcommand{\arraystretch}{1.3}
\caption{Black-box runs. Each row changes one mechanism relative to the row above.}
\label{tab:blackbox}
\centering
\begin{tabular}{l l}
\hline\hline
\textbf{Change} & \textbf{Outcome} \\
\hline
baseline                         & $\hat{p}$ pinned at $0.999$, power $4 \to 33$ \\
hard power cap, no entropy bonus & power fixed at $2.00$, every metric flat \\
top-$K$ subcarrier mask, $K=8$   & $\hat{p} \in [0.985, 0.996]$, entropy still flat \\
$K = 8 \to 1$                    & $\hat{p} \in [0.35, 0.44]$ \\
action held over the frame       & BER $\approx 0.013$, $\hat{p} \approx 0.3$ \\
\hline\hline
\end{tabular}
\end{table}

Two rows are not what they look like.
The mask is a real prior rather than a neutral change, since it assumes the answer is sparse instead of
letting the agent find that, and it is defensible only as a feasibility check.
And $K=1$ landing at $\hat{p} \approx 0.4$ instead of the $0.0002$ the probe predicted is not learning
either: the flow resamples at every OFDM symbol and may pick a different subcarrier each time, so one
hopping tone looks like many tones in a frame-level spectrogram, and holding the action over the frame
removes the gap.
The apparent threshold at four active subcarriers is a property of our sweep and not of the detector,
as Section~\ref{subsec:exphist:characterisation} shows.

The problem is the action space, not the reward.
Raw IQ gives a search space where the reward is constant everywhere the policy can reach, and no
reward shaping supplies the per-dimension credit that is missing.
